\section{Recovering the local data}
\label{local:section}

The commutator detects a relation between multiplicative characters.
That relation recovers quasi-prime divisors, including those that value
constants nontrivially.  At this stage their residue fields may have
characteristic~$\ell$.  We therefore work with character groups throughout
and use a Galois interpretation only at the top field.
For the commuting-pair approach to valuation detection and its
development via rigid elements, see
\cite{Bogomolov1991,BT2002,TopazCL2017}.
We use the character-space formulation below to include the required
residue fields.

\subsection{Commutators and alternating characters}

For any field $F$, put
\[
 \mathcal G_F=W_F\otimes_{\Zell}\Qell.
\]
A pair $f,g\in\mathcal G_F$ is \emph{alternating} if
\begin{equation}\label{local:alternating}
 f(x)g(1-x)=f(1-x)g(x)\qquad(x\in F\setminus\{0,1\}).
\end{equation}
A subspace is alternating if each pair of its elements is alternating.
Every character in $\mathcal G_F$ kills $-1$.

\begin{lemma}\label{local:commuting}
Let $F/\kappa$ be a function field over an algebraically closed field
of characteristic different from~$\ell$.  Under Kummer duality $W_F=\Pi_F^a$, two elements of $W_F$
are alternating if and only if their commutator in $\Delta_F$ is zero.
This assertion holds also for $\ell=2$.
\end{lemma}

\begin{proof}
We give the finite-coefficient argument, keeping track of the
infinite rank of $W_F$.  Write $G=G_F^{(\ell)}$, $A=G^{\mathrm{ab}}$
and $R=\mathbb Z/\ell^e$, with trivial action.  Kummer duality makes
$A$ a product of copies of~$\Zell$.

First, a class in $H^2_{\mathrm{cont}}(A,R)$ is determined by its
commutator form.  Indeed, a class with zero commutator defines
an abelian extension of $A$ by $R$.  Such an extension splits
continuously: the Pontryagin dual of $A$ is a direct sum of copies of
$\Qell/\Zell$, hence is divisible and injective.  Conversely every
continuous alternating bilinear form $A\times A\to R$ is the
commutator form of a cup-product class.  To see this, continuity and
compactness give an open subgroup in the radical of the form, so it
depends on finitely many coordinates.  In coordinate characters $x_i$
it is a sum of the forms
\[
 x_i(a)x_j(b)-x_j(a)x_i(b),\qquad i<j,
\]
which come from the cocycles $x_i\cup x_j$.
Consequently cup products generate $H^2_{\mathrm{cont}}(A,R)$.
This proof uses no division by~$2$.  In particular, even when $\ell=2$,
a diagonal cup has zero commutator and therefore vanishes.  No extra
power or Bockstein class survives on the torsion-free group~$A$.

Next, inflation
\[
 H^2_{\mathrm{cont}}(G,R)\longrightarrow H^2(G_F,R)
\]
is injective.  Let $F(\ell)$ be the maximal Galois pro-$\ell$ extension
of $F$.  An element of $F(\ell)^\times$ lies in a finite Galois
$\ell$-extension $F'/F$.  Adjoin to $F'$ the $\ell^e$-th roots of all
its $F$-conjugates.  Since the roots of unity lie in $F$, the resulting
field is again a finite Galois $\ell$-extension of $F$.  Thus
$F(\ell)^\times$ is $\ell^e$-divisible.  Kummer theory gives
$H^1(G_{F(\ell)},R)=0$, and the five-term sequence proves the claimed
injectivity.

Fix compatible identifications of the roots of unity with the
coefficient groups.  The degree-two norm-residue theorem
\cite{MS1982} identifies the kernel of the tensor cup map
\[
 (F^\times/F^{\times\ell^e})\otimes_R
 (F^\times/F^{\times\ell^e})\longrightarrow H^2(G_F,R)
\]
with the subgroup generated by the Steinberg tensors
$[x]\otimes[1-x]$.  Since the tensor cup map onto
$H^2_{\mathrm{cont}}(A,R)$ is surjective, the kernel of inflation from
$A$ to $G$ is generated by the Steinberg cups $[x]\cup[1-x]$.

Finally, put $N=[G,G]$.  The five-term sequence for $N\subset G$
identifies this same kernel with
\[
 \operatorname{Hom}_{\mathrm{cont}}(N/[N,G],R)
 =\operatorname{Hom}_{\mathrm{cont}}(\Delta_F,R).
\]
The commutator form of the transgression of a character $\chi$ is
$\chi([\ ,\ ]_F)$, up to the uniform sign determined by the
transgression convention.  Its value at $f,g$ therefore vanishes for
all~$\chi$ precisely when
\[
 f(x)g(1-x)-f(1-x)g(x)=0\pmod{\ell^e}
\]
for every $x\ne0,1$.  Continuous characters to finite cyclic
$\ell$-groups separate points of $\Delta_F$.  Letting $e$ vary proves
the assertion.
\end{proof}

It follows that a bracket-compatible isomorphism preserves alternating
pairs on $W_F$, and hence on $\mathcal G_F$: multiply rationalized
characters by powers of~$\ell$ and use
\eqref{local:alternating}.  The latter relation makes sense on residue
fields of every characteristic, without invoking their Galois groups.

\subsection{Character groups of valuations}

We regard equivalent valuations as the same valuation.  Write $v\leq w$
if $v$ is a coarsening of $w$.  For a valuation $v$ of $F$, let
$\mathcal O_v$, $\mathfrak m_v$, $Fv$ and $vF$ denote its valuation
ring, maximal ideal, residue field and value group.  Set
\[
 U_v=\mathcal O_v^\times,\qquad U_v^1=1+\mathfrak m_v,
 \qquad I_v=U_v^\perp\subset D_v=(U_v^1)^\perp\subset W_F,
\]
and write $\mathcal I_v=I_v\otimes\Qell$ and
$\mathcal D_v=D_v\otimes\Qell$.
These are the \emph{minimized} inertia and decomposition groups.
When both $F$ and $Fv$ have characteristic different from~$\ell$,
they agree with the usual abelian inertia and decomposition groups.
Their definitions above impose no characteristic restriction.

A subset of $\mathcal G_F$ is \emph{valuative} if it is contained in
$\mathcal I_v$ for some valuation~$v$.
We use the following local character theorem
\cite[Fact~2.1, Lemma~2.5 and Corollary~2.7]{Topaz2019}:
\begin{enumerate}
 \item Every valuative subset $\Sigma$ has a unique coarsest
 valuation $v_\Sigma$ with $\Sigma\subset\mathcal I_{v_\Sigma}$.
 It coarsens every valuation with this property.
 \item In an alternating subspace $S$, the valuative elements form
 a valuative subspace $\Sigma$ of codimension at most one, and
 $S\subset\mathcal D_{v_\Sigma}$.
 \item If $\Sigma$ is valuative and a character $f$ alternates with
 every member of $\Sigma$, then $f\in\mathcal D_{v_\Sigma}$.
\end{enumerate}
These statements use $\operatorname{Hom}(F^\times,\Zell)\otimes\Qell$
and hold for arbitrary fields~$F$.  The condition on the value at $-1$
in the local character theorem is automatic here.  This is the form
needed below when $\operatorname{char}(Fv)=\ell$.

Let now $F/\kappa$ be a function field over an algebraically closed
field, of transcendence degree~$s$.  A \emph{quasi-prime divisor} is a
valuation minimal among those satisfying
\[
 vF/v\kappa\cong\mathbb Z,
 \qquad \operatorname{trdeg}(Fv/\kappa v)=s-1.
\]
A quasi-prime $r$-divisor is a composition of $r$ successive
quasi-prime divisors of the successive residue function fields.  It
satisfies
\begin{equation}\label{local:rank}
 vF/v\kappa\cong\mathbb Z^r,
 \qquad \operatorname{trdeg}(Fv/\kappa v)=s-r.
\end{equation}
The residue constants $\kappa v$ are algebraically closed.

We recall the finite-generation point implicit in
\eqref{local:rank}.  The valuation transcendence-degree inequality is
\begin{equation}\label{local:abhyankar}
 \operatorname{rrank}(vF/v\kappa)
 +\operatorname{trdeg}(Fv/\kappa v)\leq s,
\end{equation}
where $\operatorname{rrank}$ is dimension after tensoring with
$\mathbb Q$.  At equality, choose relative value-independent elements
and elements with algebraically independent residues.  Together they
form a transcendence basis, and $F$ is finite over the resulting
rational field.  The finite-extension valuation inequality then shows
that the relative value group is finitely generated and the residue
extension is finitely generated.  Since $v\kappa$ is divisible,
$vF/v\kappa$ is torsion free, hence is a finitely generated free
abelian group.  Applying this argument successively gives
\eqref{local:rank}.

\begin{lemma}\label{local:residue}
For a quasi-prime $r$-divisor $v$ of $F/\kappa$, restriction to units
induces a canonical topological isomorphism
\[
 D_v/I_v\cong W_{Fv}.
\]
If $f,g\in\mathcal D_v$, they are alternating if and only if their
residual characters in $\mathcal G_{Fv}$ are alternating.
\end{lemma}

\begin{proof}
There is an exact sequence of abelian groups
\[
 0\longrightarrow (Fv)^\times/(\kappa v)^\times
 \longrightarrow F^\times/(\kappa^\times U_v^1)
 \longrightarrow vF/v\kappa\longrightarrow0.
\]
Its final term is free, so applying $\operatorname{Hom}(-,\Zell)$
gives a surjection $D_v\to W_{Fv}$ with kernel $I_v$.
This map is continuous for pointwise convergence; compactness makes
the induced bijection on the quotient a topological isomorphism.

For the second assertion, use \eqref{local:alternating}.  If $v(x)>0$,
then $1-x$ is a principal unit.  If $v(x)<0$, the identity
$1-x=-x(1-x^{-1})$ makes the characters of $1-x$ equal those of~$x$.
If $v(x)=0$ and $v(1-x)>0$, then $x$ is a principal unit.  In each
case the alternating determinant is zero.  In the remaining case
$x$ and $1-x$ are units with residues different from $0,1$, and the
determinant is exactly the residual determinant.  Conversely every
residual element different from $0,1$ has such a lift.
\end{proof}

In particular, a unit representing a class modulo $\ell^e$ evaluates
on $D_v/I_v$ exactly as its residue does on $W_{Fv}$.  This observation
will allow finite Kummer tests to pass through a valuation.

\subsection{An intrinsic description of quasi-prime divisors}

The next argument follows the character-space proof of
\cite[Fact~3.2 and Theorem~3.3]{Topaz2019}.  We include it because we
need it also in residue characteristic~$\ell$, whereas that section
of the cited paper is stated in a cohomological setting with
characteristic different from~$\ell$.

\begin{lemma}\label{local:dimension}
For $F/\kappa$ as above, in any characteristic, $s$ is the maximum
dimension of an alternating subspace of $\mathcal G_F$.
If $S$ has dimension $s$, its valuative subspace is
$\mathcal I_u$ for its associated valuation $u$, and $u$ attains
equality in \eqref{local:abhyankar}.
\end{lemma}

\begin{proof}
For every valuation $u$, characters kill the divisible constant group,
so
\[
 \dim_{\Qell}\mathcal I_u
 \leq\operatorname{rrank}(uF/u\kappa).
\]
Let $S$ be alternating and let $\Sigma$ be its valuative subspace,
with associated valuation~$u$.  If $S=\Sigma$, the asserted dimension
bound follows from \eqref{local:abhyankar}.
Otherwise $\Sigma$ has codimension one in~$S$.
A character in $S\setminus\Sigma\subset\mathcal D_u$ has nonzero
residual character.  Thus $Fu\ne\kappa u$, and
$\operatorname{trdeg}(Fu/\kappa u)\geq1$.  Consequently
\[
 \dim S\leq\operatorname{rrank}(uF/u\kappa)+1\leq s.
\]
This also bounds arbitrary alternating subspaces, by applying the
argument to their finite-dimensional subspaces.

There is a full discrete flag at a smooth closed point on a model
of $F/\kappa$.  Its inertia space has dimension~$s$ and is
alternating by the valuation inequality for $x$ and $1-x$.
Thus the maximum equals~$s$.

If $\dim S=s$ and $S=\Sigma$, all inequalities
\[
 s\leq\dim\mathcal I_u
 \leq\operatorname{rrank}(uF/u\kappa)\leq s
\]
are equalities.  If $S\ne\Sigma$, all inequalities
\[
 s-1\leq\dim\mathcal I_u
 \leq\operatorname{rrank}(uF/u\kappa)
 \leq s-\operatorname{trdeg}(Fu/\kappa u)\leq s-1
\]
are equalities.  In both cases $\Sigma=\mathcal I_u$ and $u$ has
no transcendence defect.
\end{proof}

\begin{lemma}\label{local:quasiprime}
Assume $s\geq2$.  A line $H\subset\mathcal G_F$ is the inertia
space of a quasi-prime divisor if and only if there are two
$s$-dimensional alternating subspaces $S_1,S_2$ with
$S_1\cap S_2=H$.  The valuation is unique.  If $0\ne h\in H$, its
decomposition space is
\[
 \mathcal D_v=\{f\in\mathcal G_F:(h,f)\text{ is alternating}\}.
\]
\end{lemma}

\begin{proof}
Suppose first that $v$ is quasi-prime.  The residue function field
$Fv/\kappa v$ has two independent discrete valuations of full rank
$s-1$, trivial on the constants.  For example, choose two independent
prime divisors on a rational subfield on a transcendence basis,
complete each to a flag, and prolong the flags to the finite
extension~$Fv$.  Independent valuation approximation
\cite[Theorem~1.2]{ADF2020} makes their inertia spaces disjoint.
Composing with $v$ gives two alternating inertia spaces of dimension
$s$ whose intersection is $\mathcal I_v$.

Conversely, let $S_1,S_2$ be as stated.  By
Lemma~\ref{local:dimension}, their valuative parts are
$\Sigma_i=\mathcal I_{u_i}$, where each $u_i$ has no transcendence
defect and $\dim\Sigma_i\geq s-1$.
We first show that $H$ is valuative.  If, say, $u_1\leq u_2$, then
\[
 \Sigma_1\subset\Sigma_2,
 \qquad \Sigma_1\subset S_1\cap S_2=H.
\]
Since $\dim\Sigma_1\geq1$, this proves the claim.
If the valuations are incomparable, let $w$ be their finest common
coarsening.  Applying independent valuation approximation on $Fw$
shows that the principal-unit groups of $u_1$ and $u_2$ generate
$U_w$.  Indeed, approximate a unit $a$ by $b$ for the first residual
valuation and approximate $1$ by $b$ for the second.  Then
$a=(a/b)b$ is a product of the two required principal units.
A character in $H$ annihilates both principal-unit groups,
and hence belongs to $\mathcal I_w$.  Again $H$ is valuative.
In either case,
\[
 H=\Sigma_1\cap\Sigma_2.
\]

Let $v=v_H$ be its coarsest associated valuation.  It coarsens both
$u_i$, and
\[
 H\subset\mathcal I_v
 \subset\mathcal I_{u_1}\cap\mathcal I_{u_2}=H.
\]
A coarsening of a valuation without transcendence defect also has
no transcendence defect: relative rational ranks add under
composition, and the two valuation inequalities sum to the equality
for the finer valuation.  Thus $v$ has no transcendence defect.
Its relative value group is finitely generated free, and its rank
is $\dim\mathcal I_v=1$.  Hence $vF/v\kappa\cong\mathbb Z$.
A coarsening with the same relative rank would have the same inertia
space, contradicting the defining minimality of~$v_H$.
Therefore $v$ is quasi-prime, and the coarsest-valuation description
also proves uniqueness.

Every member of $\mathcal D_v$ alternates with $h$, by the calculation
in Lemma~\ref{local:residue}.  For the converse apply the third part
of the local character theorem to the valuative line~$H$.
\end{proof}

The integral groups are recovered without an index ambiguity:
\[
 I_v=\mathcal I_v\cap W_F,
 \qquad D_v=\mathcal D_v\cap W_F.
\]
Indeed these are annihilators with values in the torsion-free group
$\Zell$.  In particular they are saturated submodules.

\begin{proposition}\label{local:transport}
Let $K/k$ and $L/l$ be function fields over algebraically closed
fields of characteristic different from~$\ell$, each of
transcendence degree at least two.  Let
$\varphi:W_L\to W_K$ be a continuous bracket-compatible
$\Zell$-module isomorphism.  Then their transcendence degrees agree,
say they are~$d$.  For every $1\leq r<d$, the map $\varphi$ carries
exactly the pairs $I_w\subset D_w$ of quasi-prime $r$-divisors of
$L/l$ onto the corresponding pairs of $K/k$.
For matched pairs it induces a continuous isomorphism
\[
 D_w/I_w\cong W_{Lw}\ \longrightarrow\
 D_v/I_v\cong W_{Kv}
\]
that preserves alternating pairs.
\end{proposition}

\begin{proof}
Lemma~\ref{local:commuting} makes the alternating relation intrinsic
to the bracket.  Lemma~\ref{local:dimension} then recovers~$d$, and
Lemma~\ref{local:quasiprime} recovers the rationalized rank-one
pairs.  Intersecting with the integral modules recovers the exact
pairs.  Lemma~\ref{local:residue} supplies the residual character
spaces and their alternating relation.  If $u$ is a valuation of $Fv$,
the inverse images of $I_u$ and $D_u$ under $D_v\to W_{Fv}$ are
$I_{u\circ v}$ and $D_{u\circ v}$, respectively, directly from the
unit and principal-unit definitions.  Apply the same
characteristic-free recognition argument to each residue function
field.  Continuing while its transcendence degree is at least two
recovers precisely the successive quasi-prime divisors through
rank~$d-1$.  The construction applies to $\varphi^{-1}$ as well.
\end{proof}

\subsection{Recognizing the finite-constant residual curves}

The remaining local input recognizes certain residual curves and all
their point inertia groups.  It does not yet distinguish true
divisors from quasi-prime divisors on the original function field.

We first describe the topological test.  A compact abelian
pro-$\ell$ group $G$ with a family $(T_i)$ of procyclic subgroups is
\emph{complete-curve-like} if one can choose generators $\tau_i$ so
that $\tau_i$ tends to zero outside finite subsets, the resulting map
\[
 \prod_i\Zell\longrightarrow G,
 \qquad(a_i)\longmapsto\sum_i a_i\tau_i
\]
has kernel the diagonal copy of $\Zell$, and its cokernel is a
finitely generated $\Zell$-module.  The diagonal kernel expresses
the single relation among the point inertias of a complete curve.
This property is invariant under continuous module isomorphisms.

For this criterion, let $F/\kappa$ be a function field over an
algebraically closed field of characteristic different from $\ell$;
its residue fields may still have characteristic $\ell$.
Form the closed subset
\[
 \mathscr I_F=
 \overline{\bigcup_{u\text{ quasi-prime}} I_u}\subset W_F.
\]
This is the closure of a \emph{union}, not the closed subgroup
generated by that union.  For a quasi-prime $(s-1)$-divisor $v$, let
\begin{equation}\label{local:residual-set}
 \mathscr R_v=
 \operatorname{im}\bigl(\mathscr I_F\cap D_v\longrightarrow
 D_v/I_v\bigr),
\end{equation}
and consider all maximal procyclic subgroups of $D_v/I_v$
contained in $\mathscr R_v$.

Pop's residual-curve criterion
\cite[Section~4, Theorem~4.2(i)(a)]{PopMin2015} applies when
\begin{enumerate}
 \item $vF$ has no nonzero $\ell$-divisible convex subgroup;
 \item there is a subfield $\kappa_1\subset F$ with
 $\kappa_1v=\kappa v$ and
 $\operatorname{trdeg}(F/\kappa_1)
 =\operatorname{trdeg}(Fv/\kappa_1v)=1$.
\end{enumerate}
It states that the group with the family just described is
complete-curve-like exactly when $\kappa v$ is an algebraic closure
of a finite field.  In that case the family is precisely all point
inertia subgroups of the smooth projective curve with function
field~$Fv$.  The criterion allows $\operatorname{char}(Fv)=\ell$;
this is the purpose of the minimized groups and of Topaz's Appendix
to~\cite{PopMin2015}.

\begin{lemma}\label{local:criterion-hypotheses}
Every quasi-prime $(s-1)$-divisor $v$ of $F/\kappa$ satisfies the two
hypotheses of the residual-curve criterion.
\end{lemma}

\begin{proof}
For a quasi-prime divisor $u$, an $\ell$-divisible convex subgroup
$H\subset uF$ has zero image in $uF/u\kappa\cong\mathbb Z$.
Thus $H\subset u\kappa$.  Coarsening by $H$ preserves the relative
rank one.  Since $u$ attains equality in the valuation inequality,
its coarsening does also, so its residual transcendence degree is
still $s-1$.  Minimality forces $H=0$.

Now consider a composite quasi-prime divisor.  The value group of
its last residual step is a nonzero convex subgroup $\Gamma$ of
the full value group and has the rank-one property just proved.
If $H$ were a nonzero $\ell$-divisible convex subgroup of the full
group, convexity would make $H\cap\Gamma$ $\ell$-divisible: an
$\ell$-th part of a positive element remains between zero and that
element.  Convex subgroups are comparable, so $H\cap\Gamma$ is
nonzero.  This contradicts the rank-one case.

For the second hypothesis choose $x_1,\ldots,x_{s-1}$ whose values
are rationally independent modulo $v\kappa$, and put
$\kappa_1=\kappa(x_1,\ldots,x_{s-1})$.
Relative value independence implies algebraic independence.  In a
polynomial in the $x_i$, distinct monomials have different values
modulo $v\kappa$, and there is a unique term of least value.
A quotient of polynomials with value zero consequently has residue
in $\kappa v$: equality of the two least values forces the same
monomial, leaving a quotient of constants.  It follows that
$\kappa_1v=\kappa v$.  The two asserted transcendence degrees are
then one by \eqref{local:rank}.
\end{proof}

\begin{proposition}\label{local:curves}
In the setting of Proposition~\ref{local:transport}, let
$w$ and $v$ be matched quasi-prime $(d-1)$-divisors.
Then $lw$ is algebraic over a finite field if and only if $kv$ is.
When this holds, the induced isomorphism $W_{Lw}\to W_{Kv}$ carries
the full family of point inertia subgroups of the respective smooth
projective curves onto each other.  Corresponding normalized order
characters differ by units of~$\Zell$.
\end{proposition}

\begin{proof}
Proposition~\ref{local:transport} recovers every group used in
\eqref{local:residual-set}.  A continuous isomorphism preserves the
closure of the union, its intersection with decomposition, its
quotient image and the maximal procyclic subgroups in that image.
Lemma~\ref{local:criterion-hypotheses} verifies the hypotheses of
Pop's criterion on both sides.  The complete-curve-like property
therefore holds simultaneously, and the criterion identifies the
entire point families when it does.  The resulting isomorphism
identifies each full procyclic subgroup with its counterpart, so it
sends a generator to a generator, differing only by a $\Zell$-unit.
\end{proof}

Only after separately excluding residue characteristic~$\ell$ will
we use these point lines to compare genera.  For the present purpose,
Proposition~\ref{local:curves} provides exact point-order vanishing
at every finite level, with no assumption on the cardinality or
transcendence degree of the original constant fields.
